\section{Introduction}
\label{sec:intro}
Reinforcement learning with verifiable rewards (RLVR), exemplified by GRPO~\citep{shao2024deepseekmath}, has emerged as a dominant paradigm for training large reasoning models~\citep{guo2025deepseek,yang2025qwen3,zeng2026glm,hong2025glm,chen2026minimax,team2026kimi}. However, GRPO derives its entire learning signal from a single outcome-level reward at the end of each rollout, which becomes prohibitively sparse in long-horizon reasoning tasks where credit must be assigned across thousands of intermediate tokens. On-policy distillation (OPD)~\citep{agarwal2024policy,song2026survey,gu2024minillm} addresses this limitation by drawing dense, token-level supervision from the logits of a stronger teacher along trajectories sampled by the student itself. The on-policy formulation mitigates the exposure bias inherent in standard supervised fine-tuning, while the refinement of supervision from the trajectory level to the token level substantially accelerates convergence. Recent studies indicate that OPD attains performance comparable to or exceeding that of RLVR, and the technique has been incorporated into the post-training pipelines of several flagship large language models, including Qwen3~\citep{yang2025qwen3}, GLM-5~\citep{zeng2026glm}, the MiMo series~\citep{xiao2026mimo}, and DeepSeek-V4~\citep{deepseekai2026deepseekv4}.

The practical applicability of OPD, however, is constrained by several requirements. The method requires white-box access to the teacher's token-level logits, precluding strong closed-source teachers; the teacher and student must share an identical vocabulary, a condition rarely satisfied across model families; and serving a separate, typically larger teacher alongside the student throughout RL training imposes substantial memory and latency overhead. On-policy self-distillation (OPSD)~\citep{zhao2026self,hubotter2026reinforcement,yang2026self,he2026self,li2026unifying,jin2026unisd,kim2026rebellious,shen2026anti,shen2026generic,ke2026respecting,wang2026trace} has recently emerged as a means of circumventing these obstacles. Under this paradigm, a single model assumes both roles, with the teacher granted access to privileged context (such as a verified reference solution) unavailable to the student. 

Existing OPSD methods~\citep{zhao2026self,hubotter2026reinforcement,yang2026self} derive token-level supervision from $e_{c,t} = \log \pi_T(y_t\mid x, y_c^*, y_{<t}) - \log \pi_S(y_t\mid x, y_{<t})$, where $x$ is the query, $y$ is a student-sampled rollout, $T$ and $S$ refer to the same model under different conditioning, and $y^*_c$ is the privileged context, typically a ground truth referenced solution provided in the dataset or a correct same-group rollout. Conditioning on a verified solution does sharpen teacher confidence on task-bearing tokens, but it also induces systematic, correctness-orthogonal shifts in generative behavior: the privileged teacher favors shorter, more assertive phrasings, suppresses exploratory hedges, and reallocates probability mass toward formatting and discourse tokens. We call this parasitic component \textbf{\emph{privilege-induced style drift}}. As shown in Figure~\ref{fig:teaser-a}, stylistic tokens dominate the learning signal, diluting the signal on task-bearing tokens and leading to the pathologies illustrated in Figure~\ref{fig:teaser-b}: \textbf{premature response-length shrinkage that limits later-stage reasoning performance, sometimes accompanied by training instability}.

% To address these issues, we propose \textsc{RLCSD}: \textbf{\underline{R}}einforcement \textbf{\underline{L}}earning with \textbf{\underline{C}}ontrastive on-policy \textbf{\underline{S}}elf-\textbf{\underline{D}}istillation. Rather than relying on $e_{c,t}$ alone as the learning signal, we form a contrastive estimate by conditioning the teacher on a correct and an incorrect solution in parallel, yielding
% \begin{gather}
% e_{c,t} = \log \pi_T(y_t \mid x, y_c^*, y_{<t}) - \log \pi_S(y_t \mid x, y_{<t}), \\
% e_{w,t} = \log \pi_T(y_t \mid x, y_w^*, y_{<t}) - \log \pi_S(y_t \mid x, y_{<t}), \\
% e_{\text{ctr},t} = e_{c,t} - e_{w,t},
% \end{gather}
% where $y_c^*$ and $y_w^*$ denote a correct and an incorrect reference solution, respectively. As standard datasets seldom provide negative references, we draw both from the student's own rollouts of the same query, taking a correct and an incorrect trajectory from the same group. The contrastive difference $e_{\text{ctr},t} = e_{c,t} - e_{w,t}$ then serves as the token-level supervision signal. Provided that the correct and incorrect hints are presented under an identical prompt template, the subtraction can suppress the stylistic component shared across the two conditionings, thereby making the remaining signal more reflective of task correctness. As illustrated in Figure~\ref{fig:teaser-a}, tokens ranked by $|e_c|$ on math problems are dominated by stylistic markers like ``Wait'' and ``Therefore'', whereas those ranked by $|e_\text{ctr}|$ more often correspond to task-bearing tokens, such as mathematical content tokens in the math reasoning task. Given the contrastive token-level signal $e_{\mathrm{ctr},t}$, we combine it with the query-level advantage $A_\mathrm{ORM}$ from GRPO by treating $e_{\mathrm{ctr},t}$ as a per-token modulation of $A_
% \mathrm{ORM}$ rather than as a substitute for it. Several further design choices, detailed in Section~\ref{sec:method}, ensure that the modulation neither reverses the direction of $A_\mathrm{ORM}$ nor allows the token-level signal to be diluted in the aggregated loss.

To address these issues, we propose \textsc{RLCSD}: \textbf{\underline{R}}einforcement \textbf{\underline{L}}earning with \textbf{\underline{C}}ontrastive on-policy \textbf{\underline{S}}elf-\textbf{\underline{D}}istillation. Rather than relying on $e_{c,t}$ alone as the learning signal, we form a contrastive estimate by conditioning the teacher on a correct and an incorrect reference solution in parallel, yielding
\begin{gather}
e_{c,t} = \log \pi_T(y_t \mid x, y_c^*, y_{<t}) - \log \pi_S(y_t \mid x, y_{<t}), \\
e_{w,t} = \log \pi_T(y_t \mid x, y_w^*, y_{<t}) - \log \pi_S(y_t \mid x, y_{<t}), \\
e_{\text{ctr},t} = e_{c,t} - e_{w,t},
\end{gather}
where $y_c^*$ and $y_w^*$ denote a correct and an incorrect reference solution, respectively. Since datasets rarely provide incorrect reference solutions, we obtain $y_w^*$ from an incorrect student rollout for the same query. The contrastive difference $e_{\text{ctr},t}$ then serves as the token-level supervision signal. When the two references are presented using the same prompt template, this subtraction can suppress stylistic components shared across the two conditioning contexts, making the remaining signal more reflective of task correctness. As illustrated in Figure~\ref{fig:teaser-a}, tokens ranked by $|e_c|$ on math problems are dominated by stylistic markers such as ``Wait'' and ``Therefore'', whereas those ranked by $|e_{\text{ctr}}|$ more often carry task-relevant information, such as mathematical content. We then use $e_{\text{ctr},t}$ to modulate the query-level advantage $A_\mathrm{ORM}$ from GRPO at each token, rather than replace it. Additional design choices, detailed in Section~\ref{sec:method}, ensure that this modulation neither reverses the direction of $A_\mathrm{ORM}$ nor allows the token-level signal to be diluted in the aggregated loss.


We evaluate RLCSD on Qwen3~\citep{yang2025qwen3} at three scales (1.7B, 4B, 8B) and Olmo-3-7B-Think~\citep{olmo2025olmo3}, with our main evaluation covering two task families: mathematical reasoning (AMC23~\citep{amc23}, AIME24~\citep{aime24}, AIME25~\citep{aime25}) and logical reasoning (Knights \& Knaves~\citep{xie2025logic}, at both in-distribution and out-of-distribution difficulty). As shown in Figure~\ref{fig:teaser-c}, RLCSD consistently outperforms GRPO and the on-policy self-distillation baselines (OPSD~\citep{zhao2026self}, SDPO~\citep{hubotter2026reinforcement}, SRPO~\citep{li2026unifying}, RLSD~\citep{yang2026self}) across the main benchmarks. Figure~\ref{fig:teaser-b} further reveals that RLCSD maintains stable training and preserves response length throughout optimization, whereas competing OPSD baselines either destabilize or suffer length shrinkage. We additionally provide experiments on agentic tasks in Appendix~\ref{app:agentic}, further demonstrating the applicability of RLCSD beyond reasoning benchmarks.

Our contributions are summarized as follows:

\begin{itemize}
\item We characterize \emph{privilege-induced style drift}, the pathology in which conditioning the teacher on a privileged solution shifts its per-token signal toward stylistic rather than correctness-bearing tokens. To address it, we propose \textbf{RLCSD}, which is designed to suppress this drift through a symmetric contrast between correct and incorrect privileged hints, and integrates the cleaned signal into RLVR as a verifier-anchored modulation of the GRPO advantage.
\item We conduct extensive experiments on Qwen3 at three scales (1.7B/4B/8B) and Olmo-3-7B-Think, showing that RLCSD \textbf{consistently outperforms} GRPO and prior on-policy self-distillation methods on mathematical and logical reasoning while maintaining stable training and response length. Additional results on agentic tasks (Appendix~\ref{app:agentic}) further demonstrate its broader applicability.
\item Extensive ablation studies show that \textbf{the contrastive principle is general}: it can be plugged into existing on-policy self-distillation methods to improve them, and our analysis extends the insight to the broader on-policy distillation setting.
\end{itemize}






